\subsection{Riesz endomorphisms}

Let E be a vector space and let u be an endomorphism of E.

The sequence \((\operatorname{Ker}(u^{n}))_{n\in\mathbf{N}}\) of vector subspaces of E is increasing; their union is a subspace of E stable under u, called the nilspace of u.

The sequence \((\mathrm{Im}(u^{n}))_{n\in\mathbf{N}}\) of vector subspaces of E is decreasing; their intersection is a subspace of E stable under u, called the conilspace of u.

Lemma 2. --- Let E be a vector space and let u be an endomorphism of E which has an index (A, V, p. 126).

If two of the following properties are satisfied, then so is the third:

\begin{enumerate}
\def\labelenumi{(\roman{enumi})}
\item
  The endomorphism u has index 0;
\item
  The nilespace N of u is finite-dimensional;
\item
  The conilespace I of u has finite codimension.
\end{enumerate}

For every integer \(n \geqslant 0\) , the endomorphism \(u^{n}\) has an index, equal to \(n \operatorname{ind}(u)\) (A, V, p. 127, lemma 2). The sequence \((\dim(\operatorname{Ker}(u^{n})))_{n \in \mathbf{N}}\) of natural integers is increasing; for it to be stationary, it is necessary and sufficient that the nilespace of u be finite-dimensional. The sequence \((\operatorname{codim}_{\operatorname{E}}(\operatorname{Im}(u^{n})))_{n \in \mathbf{N}}\) of natural integers is increasing; for it to be stationary, it is necessary and sufficient that the conilespace of u be finite-dimensional. Suppose condition (i) is satisfied; the relation

\[
\dim (\operatorname{Ker} (u ^ {n})) - \operatorname{codim} _ {\mathrm{E}} (\operatorname{Im} (u ^ {n})) = \operatorname{ind} (u ^ {n}) = n \operatorname{ind} (u) = 0
\]

then implies the equivalence of conditions (ii) and (iii). Conversely, if conditions (ii) and (iii) are satisfied, this formula implies that the sequence \((n\operatorname{ind}(u))_{n\in\mathbf{N}}\) is stationary, so u has index zero.

Lemma 3. --- Let E be a vector space and let u be an endomorphism of E which has an index. Suppose that the index of u is zero. Let N be the nilspace of u and I its conilspace.

\begin{enumerate}
\def\labelenumi{\alph{enumi})}
\item
  The space E is then the direct sum of the subspaces N and I (Weyr–Fitting decomposition, cf. A, VIII, §2, no. 2, p. 25).
\item
  The endomorphism u induces, by passing to subspaces, a nilpotent endomorphism of N and an automorphism \(u_{I}\) of I.
\item
  Let v denote the endomorphism of E which is zero on N and which coincides with \(u_{I}^{-1}\) on I. We have \(u \circ v = v \circ u\) and the endomorphism \(1_{E} - u \circ v\) is the projector with image N and kernel I.
\item
  Every endomorphism of E that commutes with u stabilizes N and I, and commutes with v.
\end{enumerate}

The endomorphism \(u\) satisfies the properties of the preceding lemma.

Let \(n \in \mathbf{N}\) such that \(\operatorname{Ker}(u^{n}) = N\) . Then \(\operatorname{Ker}(u^{m}) = \operatorname{Ker}(u^{n})\) for every integer \(m \geqslant n\) ; since \(\operatorname{ind}(u^{m}) = m \operatorname{ind}(u) = 0\) , this entails \(\operatorname{Im}(u^{m}) = \operatorname{Im}(u^{n})\) for \(m \geqslant n\) , hence \(\operatorname{I} = \operatorname{Im}(u^{n})\) . It follows from A, VIII, p. 25, prop. 2 that assertions a) and b) are satisfied.

For every \(x \in N\) , we have \(v(x) = 0\) and \(u(x) \in N\) , hence \(u(v(x)) = 0\) and \(v(u(x)) = 0\) . For every \(x \in I\) , we have \(u(v(x)) = v(u(x)) = x\) . This implies that \(u \circ v = v \circ u\) and that \(1_{E} - u \circ v\) is the projector with image N and kernel I.

Let w be an endomorphism of E that commutes with u. Let \(n \in \mathbf{N}\) such that \(N = \operatorname{Ker}(u^{n})\) and \(I = \operatorname{Im}(u^{n})\) . Let \(x \in N\) ; we have \(u^{n}(w(x)) = w(u^{n}(x)) = 0\) , hence \(w(x) \in N\) . Let \(x \in I\) ; there exists \(y \in E\) such that \(x = u^{n}(y)\) ; then, \(w(x) = w(u^{n}(y)) = u^{n}(w(y)) \in I\) . This proves that w stabilizes N and I.

Let us finally prove that \(w\) and \(v\) are permutable. For every \(x\) in \(N\) , we have \(v(w(x)) = 0 = w(v(x))\) . Let \(x \in I\) ; the elements \(v(w(x))\) and \(w(v(x))\) of \(E\) belong to \(I\) and their images under \(u\) are both equal to \(w(x)\) , hence they are equal. By linearity, this implies that \(v\) and \(w\) commute.

Remark. --- Suppose there exists an integer \(n \geqslant 0\) such that \(\operatorname{Ker}(u^{n}) = \operatorname{Ker}(u^{n+1})\) (resp. \(\operatorname{Im}(u^{n}) = \operatorname{Im}(u^{n+1})\) ); one then has \(\operatorname{Ker}(u^{m}) = \operatorname{Ker}(u^{n})\) (resp. \(\operatorname{Im}(u^{m}) = \operatorname{Im}(u^{n})\) ) for every integer \(m \geqslant n\) .

Suppose that \(u\) has index zero. For every integer \(n \geqslant 0\) , we have

\[
n \operatorname{ind} (u) = \operatorname{ind} \left(u ^ {n}\right) = \operatorname{codim} _ {\mathrm{E}} \left(\operatorname{Im} \left(u ^ {n}\right)\right) - \dim \left(\operatorname{Ker} \left(u ^ {n}\right)\right).
\]

The conditions \(\operatorname{Ker}(u^n) = \operatorname{Ker}(u^{n+1})\) and \(\operatorname{Im}(u^n) = \operatorname{Im}(u^{n+1})\) are therefore equivalent. When they are satisfied, the vector space E is the direct sum of \(\operatorname{Ker}(u^n)\) and of \(\operatorname{Im}(u^n)\) (A, VIII, p. 25, prop. 2), the space \(\operatorname{Ker}(u^n)\) is finite-dimensional, and \(u\) induces, by passage to subspaces, an automorphism of \(\operatorname{Im}(u^n)\) .

DEFINITION 2. --- Let E be a separated locally convex space. A Riesz endomorphism of E is any Fredholm endomorphism of E whose nilespace is finite-dimensional and whose conilespace is of finite codimension.

PROPOSITION 5. --- Let E be a separated locally convex space. Every Riesz endomorphism of E has index zero.

This follows from the definition and Lemma 2.

Examples. --- 1) Every automorphism of E is a Riesz endomorphism of E. If E is finite-dimensional, every element of \(\mathcal{L}(\mathrm{E})\) is a Riesz endomorphism of E.

\begin{enumerate}
\def\labelenumi{\arabic{enumi})}
\setcounter{enumi}{1}
\item
   If \((\mathrm{E}_{i})_{i\in\mathrm{I}}\) is a finite family of separated locally convex spaces and \(u_{i}\) is an element of \(\mathcal{L}(\mathrm{E}_{i})\) for all i in I, the endomorphism \(u=\bigoplus_{i\in\mathrm{I}}u_{i}\) of the locally convex space \(\bigoplus_{i\in\mathrm{I}}E_{i}\) is a Riesz endomorphism if and only if \(u_{i}\) is one for all i.
\item
   Let E be a separated locally convex space over R, and let \(\mathrm{E}_{(\mathbf{C})}\) be its complexification. Let u be an endomorphism of E. For u to be a Riesz endomorphism, it is necessary and sufficient that its complexification \(u_{(\mathbf{C})}\) be a Riesz endomorphism of \(\mathrm{E}_{(\mathbf{C})}\) .
\end{enumerate}

PROPOSITION 6. --- Let E be a separated locally convex space and let u be a Riesz endomorphism of E. Denote by N the nilspace of u and by I its conilspace.

\begin{enumerate}
\def\labelenumi{\alph{enumi})}
\item
  The subspaces N and I of E are closed and stable under u, and the space E is their topological direct sum.
\item
  The endomorphism u defines by restriction an automorphism \(u_{I}\) of I;
\item
   The vector space N is finite-dimensional, and u defines by restriction a nilpotent endomorphism \(u_{N}\) of N;
\item
  Let v be the element of \(\mathcal{L}(\mathrm{E})\) that is zero on N and coincides with \(u_{\mathrm{I}}^{-1}\) on I. It is a Riesz endomorphism of E and a quasi-inverse of u that commutes with u. The endomorphism \(1_{\mathrm{E}} - u \circ v\) of E is the projector with image N and kernel I. Every element of \(\mathcal{L}(\mathrm{E})\) which commutes with u stabilizes N and I, and commutes with v.
\end{enumerate}

Let n be a natural number such that \(\operatorname{Ker}(u^{n}) = \operatorname{N}\) and \(\operatorname{Im}(u^{n}) = \operatorname{I}\) . Since E is separated and u is continuous, \(\operatorname{N} = \operatorname{Ker}(u^{n})\) is closed in E. Since u is a Fredholm endomorphism, so is \(u^{n}\) , and \(\operatorname{I} = \operatorname{Im}(u^{n})\) is closed (III, p. 42, prop. 2). By lemma 2, the space E is the direct sum of N and I, and assertion a) then follows from prop. 3 of EVT, I, p. 15.

Since I is closed, of finite codimension in E, and stable under u, the map \(u_{I}\colon I \to I\) induced by u is a Fredholm endomorphism (III, p. 44, prop. 3); it is also bijective, hence it is an automorphism of I (III, p. 42, prop. 2). This proves b), and the remaining assertions follow from lemma 2.

The endomorphism \(v\) defined in assertion \(d\) ) of the proposition is a quasi-inverse of \(u\) , which is called the canonical quasi-inverse of \(u\) .

Riesz endomorphisms enjoy certain of the stability properties of Fredholm maps.

PROPOSITION 7. — Let E be a separated locally convex space and let u be an endomorphism of E.

\begin{enumerate}
\def\labelenumi{\alph{enumi})}
\item
   Let \(E_{1}\) be a closed subspace of finite codimension of E, stable under u. Let \(u_{1}\) be the element of \(\mathcal{L}(E_{1})\) which coincides with u on \(E_{1}\) . Then u is a Riesz endomorphism if and only if \(u_{1}\) is one;
\item
   Suppose that u is a Fredholm endomorphism of index 0. Let \(\widehat{\mathbf{E}}\) be the completion of E, and let \(\widehat{u} \in \mathcal{L}(\widehat{\mathbf{E}})\) be the continuous extension of u. Then u is a Riesz endomorphism if and only if \(\widehat{u}\) is one;
\item
 Endow the dual E' of E with the topology of bounded convergence, the topology of compact convergence, or the weak topology. If u is a Riesz endomorphism of E, then \(^{t}\) u is a Riesz endomorphism of E'.
\end{enumerate}
a) For \(u\) to be a Fredholm endomorphism of index zero, it is necessary and sufficient that \(u_1\) be one (III, p. 44, prop. 3). On the other hand, for every \(n \in \mathbf{N}\) , we have \(\mathrm{Ker}(u_1^n) = \mathrm{E}_1 \cap \mathrm{Ker}(u^n)\) , whence


\[
\dim \operatorname{Ker} (u _ {1} ^ {n}) \leqslant \dim \operatorname{Ker} (u ^ {n}) \leqslant \dim \operatorname{Ker} (u _ {1} ^ {n}) + \operatorname{codim} _ {\mathrm{E}} (\mathrm{E} _ {1}),
\]

therefore the sequence \((\dim\operatorname{Ker}(u^{n}))_{n\in\mathbf{N}}\) is bounded if and only if the sequence \((\dim\operatorname{Ker}(u_{1}^{n}))_{n\in\mathbf{N}}\) is. Consequently, u is a Riesz endomorphism of E if and only if \(u_{1}\) is a Riesz endomorphism of \(E_{1}\) (lemma 2).

\begin{enumerate}
\def\labelenumi{\alph{enumi})}
\setcounter{enumi}{1}
\item
  By prop. 4 of III, p. 44, for every \(n \in \mathbf{N}\) , the map \(\widehat{u}^n\) is a Fredholm endomorphism of \(\widehat{\mathrm{E}}\) such that \(\operatorname{Ker}(\widehat{u}^n) = \operatorname{Ker}(u^n)\) and \(\operatorname{ind}(\widehat{u}^n) = \operatorname{ind}(u^n) = n\operatorname{ind}(u)\) . Consequently, \(u\) is a Riesz endomorphism if and only if \(\widehat{u}\) is a Riesz endomorphism.
\item
  The transpose \({}^{t}u\) is a Fredholm endomorphism of index 0 of E' (III, p. 43, no. 3). Since the kernel of \((^{t}u)^{n}\) is the orthogonal complement of the image of \(u^{n}\) (EVT, IV, p. 27, prop. 2), the sequence \((\mathrm{Ker}({}^{t}u)^{n})_{n}\) is stationary. This proves that \({}^{t}u\) is a Riesz endomorphism of E'.
\end{enumerate}

PROPOSITION 8. --- Let E be a separated locally convex space and u an endomorphism of E. The following conditions are equivalent:

\begin{enumerate}
\def\labelenumi{(\roman{enumi})}
\item
  u is a Riesz endomorphism of E;
\item
  There exists a quasi-inverse v of u that commutes with u;
\item
  There exists a closed vector subspace \(E_{1}\) of E, of finite codimension, stable under u, such that u induces an automorphism of \(E_{1}\) .
\end{enumerate}

Implication (i) \(\Longrightarrow\) (ii) follows from prop. 6, d).

Suppose that \(u\) has a quasi-inverse \(v\) which commutes with \(u\) . The kernel \(\mathrm{E}_1\) of \(1_{\mathrm{E}} - u \circ v\) is a closed vector subspace of E, of finite codimension, stable under \(u\) and \(v\) . Since one has \(u(v(x)) = v(u(x)) = x\) for all \(x \in \mathrm{E}_1\) , the maps \(u\) and \(v\) induce automorphisms of \(\mathrm{E}_1\) inverses of each other. Hence (ii) implies (iii).

Finally, the implication (iii) \(\Longrightarrow\) (i) follows from the assertion \(a\) ) of Prop. 7 and from the fact that an automorphism is a Riesz endomorphism.

PROPOSITION 9. --- Let E be a separated locally convex space, and let u, v be permutable elements of \(\mathcal{L}(\mathrm{E})\) . The following conditions are equivalent:

\begin{enumerate}
\def\labelenumi{(\roman{enumi})}
\item
  The endomorphisms u and v are Riesz endomorphisms.
\item
  The endomorphism \(u \circ v\) is a Riesz endomorphism.
\end{enumerate}

Suppose that u and v are Riesz endomorphisms; denote by \(u'\) and \(v'\) their canonical quasi-inverses. The endomorphisms u, v, \(u'\) and \(v'\) of E commute (prop. 6, d)). The endomorphism \(v' \circ u'\) is a quasi-inverse of \(u \circ v\) which commutes with \(u \circ v\) , hence \(u \circ v\) is a Riesz endomorphism (prop. 8).

Suppose conversely that \(u \circ v\) is a Riesz endomorphism, and let w be its canonical quasi-inverse. Since u commutes with \(u \circ v\) , the endomorphisms u and w commute by prop. 6, d). Likewise v and w commute. Hence the endomorphisms u, v and w of E commute; it follows that the endomorphism \(v \circ w\) of E is a quasi-inverse of u and that \(w \circ u\) is a quasi-inverse of v. By prop. 8, u and v are Riesz endomorphisms.

\pandocbounded{\includegraphics[keepaspectratio]{images/part001_6cebf9a3.jpg}}

Let u be a Riesz endomorphism of E. If v is an automorphism of E, the endomorphism \(u \circ v\) of E is not necessarily a Riesz endomorphism of E, even if \(u \circ v - v \circ u\) has finite rank. If h is a finite-rank endomorphism of E, the endomorphism \(u + h\) is not necessarily a Riesz endomorphism (cf. III, p. 120, exercise 3).

PROPOSITION 10. --- Let E and F be locally convex spaces, \(p \in \mathcal{L}(\mathrm{E};\mathrm{F})\) and \(q \in \mathcal{L}(\mathrm{F};\mathrm{E})\) . Set \(u = 1_{\mathrm{E}} - q \circ p\) and \(v = 1_{\mathrm{F}} - p \circ q\) .

\begin{enumerate}
\def\labelenumi{\alph{enumi})}
\item
  Let \(n \in \mathbf{N}\) . The map \(p\) defines by restriction an isomorphism of topological vector spaces from \(\mathrm{Ker}(u^n)\) onto \(\mathrm{Ker}(v^n)\) and defines by passage to quotients an isomorphism of topological vector spaces from Coker( \(u^n\) ) onto Coker( \(v^n\) );
\item
  If the image of u is closed in E, that of v is closed in F;
\item
   If \(\operatorname{Ker}(u)\) has a topological supplement in E, then \(\operatorname{Ker}(v)\) has one in F. If \(\operatorname{Im}(u)\) has a topological supplement in E, then \(\operatorname{Im}(v)\) has one in F;
\item
  If u is a strict morphism, then v is a strict morphism.
\item
  If u is an automorphism of E, then v is an automorphism of F.
\item
  If \(u\) is a Fredholm endomorphism of E, then \(v\) is a Fredholm endomorphism of F, and we have \(\operatorname{ind}(u) = \operatorname{ind}(v)\) ;
\item
  Suppose E and F are separated. If u is a Riesz endomorphism of E, then v is a Riesz endomorphism of F.
\end{enumerate}

Let us first note the formulas

\[
q \circ v = u \circ q, \quad v \circ p = p \circ u.\tag{6}
\]

Let us then prove two lemmas.

Lemma 4. --- Let \(n \in \mathbf{N}\) . Set

\[
q _ {n} = \sum_ {i = 1} ^ {n} (- 1) ^ {i - 1} \binom {n} {i} q \circ (p \circ q) ^ {i - 1}.
\]

One has \(u^{n}=1_{E}-q_{n}\circ p\) and \(v^{n}=1_{F}-p\circ q_{n}\) .

One has \((q\circ p)^{i} = q\circ (p\circ q)^{i - 1}\circ p\) for every integer \(i\geqslant 1\) . We then compute

\[
u ^ {n} = (1 _ {\mathrm{E}} - q \circ p) ^ {n} = 1 _ {\mathrm{E}} - \sum_ {i = 1} ^ {n} (- 1) ^ {i - 1} {\binom {n} {i}} (q \circ p) ^ {i} = 1 _ {\mathrm{E}} - q _ {n} \circ p
\]

\[
v ^ {n} = (1 _ {\mathrm{F}} - p \circ q) ^ {n} = 1 _ {\mathrm{F}} - \sum_ {i = 1} ^ {n} (- 1) ^ {i - 1} \binom {n} {i} (p \circ q) ^ {i} = 1 _ {\mathrm{F}} - p \circ q _ {n},
\]

whence the result.

Lemma 5. --- Let \(u'\) be an element of \(\mathcal{L}(\mathrm{E})\) . Set \(v' = 1_{\mathrm{F}} + p \circ u' \circ q\) . If \(u'\) is an inverse of \(u\) , then \(v'\) is an inverse of \(v\) . If \(u'\) is a quasi-inverse of \(u\) , then \(v'\) is a quasi-inverse of \(v\) . If \(u'\) commutes with \(u\) , then \(v'\) commutes with \(v\) .

Using formulas (6), we compute

\[
v - v ^ {\prime} \circ v = - p \circ u ^ {\prime} \circ q \circ v = - p \circ u ^ {\prime} \circ u \circ q
\]

\[
v - v \circ v ^ {\prime} = - v \circ p \circ u ^ {\prime} \circ q = - p \circ u \circ u ^ {\prime} \circ q,
\]

whence

\[
1 _ {\mathrm{F}} - v ^ {\prime} \circ v = p \circ (1 _ {\mathrm{E}} - u ^ {\prime} \circ u) \circ q
\]

\[
1 _ {\mathrm{F}} - v \circ v ^ {\prime} = p \circ (1 _ {\mathrm{E}} - u \circ u ^ {\prime}) \circ q,
\]

which proves the lemma.

Given Lemma 4, it suffices to prove assertion \(a\) ) for \(n = 1\) . According to formulas (6), the maps \(p\) and \(q\) define, by passing to subspaces, continuous linear maps \(p'\) : \(\operatorname{Ker}(u) \to \operatorname{Ker}(v)\) and \(q'\) : \(\operatorname{Ker}(v) \to \operatorname{Ker}(u)\) , and by passing to quotients of continuous linear maps \(p''\) : \(\operatorname{Coker}(u) \to \operatorname{Coker}(v)\) and \(q''\) : \(\operatorname{Coker}(v) \to \operatorname{Coker}(u)\) .

Since \(1_{\mathrm{E}} - u = q \circ p\) and \(1_{\mathrm{F}} - v = p \circ q\) , the maps \(p'\) and \(q'\) on the one hand, \(p''\) and \(q''\) on the other hand, are mutually inverse isomorphisms. This proves \(a\) ).

For the image of a continuous linear map to be closed, it is necessary and sufficient that the cokernel of this map be a separated space. Thus b) follows from a).

Denote by \(i\colon \operatorname{Ker}(u)\to \mathrm{E}\) and \(j\colon \operatorname{Ker}(v)\to \mathrm{F}\) the canonical injections. Suppose that \(\operatorname{Ker}(u)\) has a topological supplement in E, and let \(r\colon \mathrm{E}\to \operatorname{Ker}(u)\) be a continuous linear retraction of \(i\) . Set \(r' = p'\circ r\circ q\) . We have \(r'\in \mathcal{L}(\mathrm{F};\operatorname {Ker}(v))\) and

\[
r ^ {\prime} \circ j = p ^ {\prime} \circ r \circ q \circ j = p ^ {\prime} \circ r \circ i \circ q ^ {\prime} = p ^ {\prime} \circ q ^ {\prime} = 1 _ {\operatorname{Ker} (v)},
\]

hence \(r'\) is a continuous linear retraction of j and \(\operatorname{Ker}(v)\) has a topological supplement in F. This proves the first assertion of c). The second is proved analogously, noting that if \(s: \operatorname{Coker}(u) \to \operatorname{E}\) is a continuous linear section of the canonical surjection from E onto \(\operatorname{Coker}(u)\) , then \(p \circ s \circ q''\) is a continuous linear section of the canonical surjection of F onto \(\operatorname{Coker}(v)\) .

Denote by \(\overline{u}\colon\mathrm{E}/\mathrm{Ker}(u)\to\mathrm{Im}(u)\) and \(\overline{v}\colon\mathrm{F}/\mathrm{Ker}(v)\to\mathrm{Im}(v)\) the continuous linear bijective mappings deduced from u and v. By formulas (6), the map p defines, by passing to quotients, a continuous linear map \(p_{1}\colon\mathrm{E}/\mathrm{Ker}(u)\to\mathrm{F}/\mathrm{Ker}(v)\) and p, q define, by passing to subspaces, continuous linear maps \(p_{0}\colon\mathrm{Im}(u)\to\mathrm{Im}(v)\) and \(q_{0}\colon\mathrm{Im}(v)\to\mathrm{Im}(u)\) .

Denote by \(t \colon \operatorname{Im}(v) \to \operatorname{F} / \operatorname{Ker}(v)\) the composite map of the canonical injection \(\operatorname{Im}(v) \to \operatorname{F}\) and of the canonical surjection \(\operatorname{F} \to \operatorname{F} / \operatorname{Ker}(v)\) ; it is continuous. We have \(\overline{v} \circ t = 1_{\operatorname{Im}(v)} - p_0 \circ q_0\) and \(\overline{v} \circ p_1 = p_0 \circ \overline{u}\) , whence

\[
\overline {{v}} \circ (t + p _ {1} \circ \overline {{u}} ^ {- 1} \circ q _ {0}) = \overline {{v}} \circ t + p _ {0} \circ q _ {0} = 1 _ {\mathrm{Im} (v)}.
\]

This proves that \(t + p_{1} \circ \overline{u}^{-1} \circ q_{0}\) is the inverse bijection of \(\overline{v}\) . If u is strict, the map \(\overline{u}^{-1}\) is continuous, and so is \(\overline{v}^{-1}\) which proves that v is strict. This proves d).

By Lemma 5, if \(u\) is an automorphism of E, then \(v\) is an automorphism of F. Similarly, if \(u\) is a Fredholm endomorphism of E, then \(v\) is a Fredholm endomorphism of F, and one then has \(\operatorname{ind}(u) = \operatorname{ind}(v)\) according to \(a\) ). This proves \(e\) ) and \(f\) ). Finally, if E and F are separated and \(u\) is a Riesz endomorphism of E, then \(v\) is a Riesz endomorphism of F by Lemma 5 and condition (ii) of Prop. 8; this is assertion \(g\) ).

